La figure ci-dessous représente une scie circulaire de rayon $R$ qui peut
tourner autour de son axe. Une courroie liée la poulie ($P_{1}$) d'un moteur
électrique et la poulie ($P_{2}$) de la scie.

La courroie ne glisse pas sur les deux poulies. L'arbre du moteur effectue
$\qty{1800}{tr\per\min}$.

\begin{figure}[H]
    \centering
    \begin{tikzpicture}[>=Stealth]
        \newdimen\ri \ri=6mm
        \newdimen\rii \rii=10mm
        \newdimen\R \R=20mm
        \draw[decoration={zigzag,amplitude=1mm,segment length=4.163mm},
            decorate,fill=gray] (0,0) circle (1.06\R);
        \node[circle,draw,fill=lightgray,inner sep=0pt,
            minimum size=2\R] (S) {};
        \foreach \x/\r [count=\i] in {-6/\ri,0/\rii}
        \node[circle,draw,fill=gray,inner sep=0pt,minimum size=2\r,
            label=below:$(P_{\i})$] (P\i) at (\x,0) {};
        \draw[thick] (P2.north) arc (90:-90:\rii) --
            (P1.south) arc (-90:-270:\ri) --
            node[above,sloped,pos=0.4] {courroie} (P2.north);
        \foreach \i in {1,2}{
            \node[circle,fill,inner sep=0.6pt] at (P\i.center) {};
            \draw[<->,shorten <=1.2pt,shorten >=1pt]
                (P\i.center) -- (P\i.-45)
                node[pos=0.5,above right=-1mm,font=\scriptsize] {$r_{\i}$};
        }
        \draw[<->,shorten <=1.2pt,shorten >=1pt]
            (P2.center) -- (S.45)
            node[pos=0.6,below right=-1mm,font=\small] {$R$};
        \node[above=3mm] at (P1.north) {moteur};
        \node[above=3mm] at (S.north) {scie};
    \end{tikzpicture}
    \caption{}
    \label{fig:scie}
\end{figure}

\begin{enumerate}
    \item Calculer la vitesse angulaire de l'arbre du moteur. 
    \item Déterminer la vitesse linéaire d'un point de la courroie. 
    \item En déduire la fréquence de rotation de la scie. 
    \item Trouver la vitesse d'une des dents de la scie. 
\end{enumerate}

\textbf{Données~:} Rayons des poulies ($P_{1}$) et ($P_{2}$) sont~:
$r_{1}=\qty{10}{\cm}$, $r_{2}=\qty{20}{\cm}$, $R=\qty{40}{\cm}$.

